\chapter{Fundamentação Teórica}
\label{sec:fundamentacao}

Neste capítulo, são apresentados os principais conceitos que embasam o desenvolvimento de \textit{pipelines} de dados com ferramentas \acs{open_source}. A seguir, aborda-se a evolução das abordagens de \ac{etl}, os desafios contemporâneos da \ac{ed} e os benefícios da adoção de soluções de código aberto.

\section{ETL e a Evolução das Arquiteturas de Dados}
\label{sec:etl_evolucao}

O processo de \acs{etl} consolidou-se historicamente como o componente fundamental na \acs{ed}, sendo responsável por extrair informações de diversas fontes, transformá-las conforme necessidades analíticas e carregá-las em ambientes de destino, como o \acs{dw} ou o \acs{dl} \cite{septian2019, nwokeji2021}. Enquanto o \acs{dw} atua como uma estrutura centralizada e desnormalizada, projetada para o suporte à decisão estratégica via análise histórica e operações \ac{olap}, o \acs{dl} surgiu como uma abordagem eficiente e escalável para o armazenamento de grandes volumes de dados brutos, permitindo o processamento sob demanda com suporte a ferramentas \acs{open_source}, como o \textit{Hadoop File System}, garantindo baixo custo, tolerância a falhas e a flexibilidade de esquema (\textit{schema on read}) \cite{septian2019}.

Essa evolução arquitetural é um reflexo direto dos desafios impostos pelo \textit{Big Data}. Segundo \cite{nwokeji2021}, o \acs{etl} evoluiu para endereçar os cinco pilares fundamentais do \textit{Big Data}, volume, variedade, velocidade, veracidade e valor. Contudo, os processos tradicionais de \acs{etl}, tipicamente confinados ao contexto de \ac{bi}, frequentemente demonstram limitações em ambientes modernos, especialmente no que tange ao custo elevado e à baixa capacidade de ingestão em tempo real \cite{nwokeji2021}. Para superar tais restrições, observou-se uma diversificação nas abordagens: \cite{boulahia2024} categorizam o \acs{etl} atual em três vertentes: o \acs{etl} tradicional para \ac{bi} (focado em dados estruturados), o \acs{etl} para \textit{Big Data} (utilizando \textit{frameworks} distribuídos como Spark para escalabilidade) e o \acs{etl} Semântico (que aplica ontologias para integrar dados heterogêneos).

Nesse cenário, a própria lógica de processamento foi redefinida, impulsionando a transição do modelo \acs{etl} para o \ac{elt}. Esta mudança é motivada em grande parte pela crescente complexidade do \textit{Data Wrangling} — o processo de obter, organizar, compreender, extrair e formatar dados. No modelo \ac{elt}, os dados brutos são primeiramente carregados no \acs{dl} ou \acs{dw}, e a transformação ocorre posteriormente, aproveitando a robusta capacidade de processamento distribuído da plataforma de destino. Esta flexibilidade é crucial para o desenvolvimento de \textit{pipelines} híbridos e escaláveis. Apesar dos desafios inerentes ao \textit{setup} e à manutenção, a adoção de ferramentas \acs{open_source}, exemplificada por plataformas como o Apache Airflow, consolidou-se como a escolha preferencial no setor. Tais ferramentas atuam como um "contêiner aberto" para diversas soluções de \textit{data wrangling}, oferecendo a escalabilidade e a agilidade necessárias para lidar com os fluxos de dados complexos que caracterizam a atual infraestrutura de dados moderna.

\section{Open Source: Estratégia, Flexibilidade e Comparação}
\label{sec:open_source}

O uso de ferramentas \acs{open_source} vem se consolidando como alternativa viável e estratégica às soluções proprietárias em projetos de engenharia de dados. \cite[]{septian2019} demonstram que, ao construir \textit{pipelines} com softwares livres, é possível reduzir custos, adaptar a arquitetura às necessidades específicas da organização e fomentar a inovação por meio da colaboração comunitária.

Especificamente para pequenas e médias empresas, \cite[]{septian2019} apontam que soluções baseadas em software livre oferecem alternativas de baixo custo para acessar tecnologias antes restritas a grandes corporações. Ferramentas como Apache Hadoop, Spark e Kafka são amplamente utilizadas em diversos ambientes corporativos, viabilizando a escalabilidade e a interoperabilidade entre diferentes camadas da arquitetura.

A escolha entre as abordagens, no entanto, depende do contexto organizacional. \cite[]{boulahia2024} apresentam uma análise comparativa:
\begin{itemize}
\item Ferramentas Proprietárias: Oferecem maior robustez, ou seja, a capacidade de um sistema operar de forma estável e confiável mesmo sob condições adversas ou não planejadas, suporte técnico especializado e garantias contratuais.suporte técnico especializado e garantias contratuais.
\item Ferramentas \acs{open_source}: Destacam-se pela flexibilidade, personalização e baixo custo.
\end{itemize}
Desta forma, ambientes acadêmicos e empresas com foco em inovação e adaptabilidade podem beneficiar-se das soluções open source, enquanto grandes corporações, devido à necessidade de uma governança de dados robusta, que assegure a linhagem, a conformidade, a segurança e a qualidade das informações em conformidade com requisitos regulatórios rigorosos, podem ainda optar por ferramentas proprietárias (\cite[]{boulahia2024}, \cite[]{septian2019}).

\section{Técnicas Avançadas de Ingestão e Processamento}
\label{sec:avancos_tecnicas}

Com o aumento exponencial da produção de dados provenientes de dispositivos \textit{IoT}, redes sociais e sistemas transacionais, surgem novos desafios na ingestão em tempo real. \cite[]{murarka2024} enfatizam que tecnologias como Apache Kafka, Flume e NiFi desempenham papel central nesse contexto, oferecendo suporte eficiente tanto a fluxos \textit{batch} quanto \textit{streaming}, sendo alternativas ao modelo tradicional do Hadoop MapReduce por oferecerem processamento em memória e análises em tempo real.

A capacidade de lidar com diferentes padrões de fluxo de dados é crucial. Plataformas como Apache Spark, Flink  e Storm permitem a construção de \textit{pipelines} mais eficientes, com menor latência e maior resiliência.

Um aspecto relevante, e tendência crescente, é a incorporação de aprendizado de máquina (\textit{machine learning}) nos \textit{pipelines} de dados. O objetivo é automatizar etapas como detecção de anomalias, classificação e priorização de fluxos. Essa integração aumenta a eficiência operacional e abre espaço para \textit{pipelines} mais autônomos e adaptativos (\cite[]{murarka2024}).

\section{Tendências Emergentes: Reverse ETL}
\label{sec:reverse_etl}

Tradicionalmente, a arquitetura analítica de dados foca na consolidação de dados em repositórios centrais. No entanto, o conceito de \textit{Reverse ETL} (\textit{Extract, Transform, Load} Invertido) tem ganhado destaque ao propor o fluxo de dados no sentido oposto: devolver os dados já tratados, enriquecidos e analisados a sistemas operacionais e de negócio.

\cite[]{reverseetl2022} explicam que essa abordagem resolve o chamado “último quilômetro da análise de dados”, reduzindo silos de informação e permitindo que \textit{insights} estratégicos sejam consumidos diretamente em ferramentas operacionais (como \texttt{CRMs}, ferramentas de marketing ou \texttt{ERPs}). Isso alinha decisões analíticas com ações operacionais em tempo quase real.

Apesar de ser uma tendência poderosa para a \textit{Operational Analytics}, o \textit{Reverse ETL} ainda enfrenta desafios técnicos como a sincronização eficiente de dados, a integração com múltiplas \texttt{APIs} de terceiros e a gestão de custos associados à latência, temas que continuam a impulsionar o desenvolvimento de novas ferramentas no ecossistema de dados (\cite[]{reverseetl2022}).


